% GENERATED FILE. Edit canonical lessons, then run npm run generate:books.
% canonical-source-hash: fnv1a64:3faa0a85d698d280
% canonical-lessons: LA-C551-capsa, LA-C551-scrinium, LA-C551-membrana, LA-C551-pugillares, LA-C551-rescribo

\chapter{Book boxes and tablets}
\label{ch:la-book-boxes-and-tablets}
\hlchaptermodality{551}

\begin{chapteropening}
\textbf{By the end of this chapter:} \emph{I can talk about writing materials and replying to letters.}

Complete the last lesson of chapter 551: I can talk about writing materials and replying to letters.
\end{chapteropening}

\section[capsa, capsae]{capsa, capsae — a book box}

\label{lesson:LA-C551-capsa}



\begin{quote}
Before the new one: say the Latin for to read through, then the Latin for to unroll.
\end{quote}

\subsection*{You'll want to know: capsa, capsae}
\textbf{capsa, capsae} — \textquotedblleft{}a book box\textquotedblright{}.

\begin{quote}
Puer librōs in capsā portat.
\end{quote}

\begin{quote}
\emph{The boy carries his books in a book box.}
\end{quote}

Feminine, first declension, from \emph{capiō} — to hold. English \emph{case}, a box, comes from it through French.

\begin{cousinweb}
\textbf{capsa} is built on \emph{capiō} \textquotedblleft{}to take, hold\textquotedblright{} (Lewis \& Short; EtymDB). English \emph{case} \textquotedblleft{}box\textquotedblright{} comes from it through French, and \emph{capsule} from its diminutive.
\end{cousinweb}

\begin{grammarlens}[title={What you've built}]
One more word for telling, explaining and describing at B1.
\end{grammarlens}

\subsection*{Your turn}
\begin{itemize}
\raggedright
  \item \emph{Say it:} \emph{capsa, capsae}
  \item \emph{Say it:} \emph{capsa, capsae}, once more
  \item \emph{Say it:} a sentence of your own with \emph{capsa, capsae}
  \item \emph{Recall:} say the Latin for to read through, then the Latin for to unroll, then say \emph{capsa, capsae} again
\end{itemize}

\begin{tcolorbox}[breakable,colback=teal!4,colframe=teal!35!black,title={Before you move on}]
What does capsa, capsae mean? (\textquotedblleft{}A book box\textquotedblright{}.) Say it once more.
\end{tcolorbox}
\section[scrīnium, scrīniī]{scrīnium, scrīniī — a writing case, chest for papers}

\label{lesson:LA-C551-scrinium}



\begin{quote}
Before the new one: say the Latin for to unroll, then the Latin for a book box.
\end{quote}

\subsection*{You'll want to know: scrīnium, scrīniī}
\textbf{scrīnium, scrīniī} — \textquotedblleft{}a writing case, chest for papers\textquotedblright{}.

\begin{quote}
Epistulās in scrīniō servō.
\end{quote}

\begin{quote}
\emph{I keep the letters in a writing case.}
\end{quote}

Neuter, second declension. A round box for book-rolls and letters.

\begin{cousinweb}
\textbf{scrīnium} is a native Latin noun; Lewis \& Short links it with \emph{scrībō} \textquotedblleft{}to write\textquotedblright{}, but EtymDB traces it to a root meaning \textquotedblleft{}to cut\textquotedblright{}, so the origin is not settled. English \emph{shrine} comes from it.
\end{cousinweb}

\begin{grammarlens}[title={What you've built}]
One more word for telling, explaining and describing at B1.
\end{grammarlens}

\subsection*{Your turn}
\begin{itemize}
\raggedright
  \item \emph{Say it:} \emph{scrīnium, scrīniī}
  \item \emph{Say it:} \emph{scrīnium, scrīniī}, once more
  \item \emph{Say it:} a sentence of your own with \emph{scrīnium, scrīniī}
  \item \emph{Recall:} say the Latin for to unroll, then the Latin for a book box, then say \emph{scrīnium, scrīniī} again
\end{itemize}

\begin{tcolorbox}[breakable,colback=teal!4,colframe=teal!35!black,title={Before you move on}]
What does scrīnium, scrīniī mean? (\textquotedblleft{}A writing case, chest for papers\textquotedblright{}.) Say it once more.
\end{tcolorbox}
\section[membrāna, membrānae]{membrāna, membrānae — parchment; skin}

\label{lesson:LA-C551-membrana}



\begin{quote}
Before the new one: say the Latin for a book box, then the Latin for a writing case.
\end{quote}

\subsection*{You'll want to know: membrāna, membrānae}
\textbf{membrāna, membrānae} — \textquotedblleft{}parchment; skin\textquotedblright{}.

\begin{quote}
Poēta versūs in membrānā scrībit.
\end{quote}

\begin{quote}
\emph{The poet writes the verses on parchment.}
\end{quote}

Feminine, first declension, from \emph{membrum}. Thin skin, then parchment for writing. English \emph{membrane}.

\begin{cousinweb}
\textbf{membrāna} is built on \emph{membrum} \textquotedblleft{}limb\textquotedblright{} (Lewis \& Short; EtymDB): the skin over the limbs, then parchment. English \emph{membrane} comes from it.
\end{cousinweb}

\begin{grammarlens}[title={What you've built}]
One more word for telling, explaining and describing at B1.
\end{grammarlens}

\subsection*{Your turn}
\begin{itemize}
\raggedright
  \item \emph{Say it:} \emph{membrāna, membrānae}
  \item \emph{Say it:} \emph{membrāna, membrānae}, once more
  \item \emph{Say it:} a sentence of your own with \emph{membrāna, membrānae}
  \item \emph{Recall:} say the Latin for a book box, then the Latin for a writing case, then say \emph{membrāna, membrānae} again
\end{itemize}

\begin{tcolorbox}[breakable,colback=teal!4,colframe=teal!35!black,title={Before you move on}]
What does membrāna, membrānae mean? (\textquotedblleft{}Parchment; skin\textquotedblright{}.) Say it once more.
\end{tcolorbox}
\section[pugillārēs, pugillārium]{pugillārēs, pugillārium — writing tablets}

\label{lesson:LA-C551-pugillares}



\begin{quote}
Before the new one: say the Latin for a writing case, then the Latin for parchment.
\end{quote}

\subsection*{You'll want to know: pugillārēs, pugillārium}
\textbf{pugillārēs, pugillārium} — \textquotedblleft{}writing tablets\textquotedblright{}.

\begin{quote}
In pugillāribus pauca verba scrībit.
\end{quote}

\begin{quote}
\emph{He writes a few words on his writing tablets.}
\end{quote}

Masculine, third declension, plural. Small wax tablets that fitted in the hand, from \emph{pugillus} — a little fist.

\begin{cousinweb}
\textbf{pugillārēs} is built on \emph{pugillus} \textquotedblleft{}a little fist, handful\textquotedblright{}, from \emph{pugnus} \textquotedblleft{}fist\textquotedblright{} (Lewis \& Short): tablets held in the hand.
\end{cousinweb}

\begin{grammarlens}[title={What you've built}]
One more word for telling, explaining and describing at B1.
\end{grammarlens}

\subsection*{Your turn}
\begin{itemize}
\raggedright
  \item \emph{Say it:} \emph{pugillārēs, pugillārium}
  \item \emph{Say it:} \emph{pugillārēs, pugillārium}, once more
  \item \emph{Say it:} a sentence of your own with \emph{pugillārēs, pugillārium}
  \item \emph{Recall:} say the Latin for a writing case, then the Latin for parchment, then say \emph{pugillārēs, pugillārium} again
\end{itemize}

\begin{tcolorbox}[breakable,colback=teal!4,colframe=teal!35!black,title={Before you move on}]
What does pugillārēs, pugillārium mean? (\textquotedblleft{}Writing tablets\textquotedblright{}.) Say it once more.
\end{tcolorbox}
\section[rescrībō, rescrībere]{rescrībō, rescrībere — to write back, reply}

\label{lesson:LA-C551-rescribo}



\begin{quote}
Before the new one: say the Latin for parchment, then the Latin for writing tablets.
\end{quote}

\subsection*{You'll want to know: rescrībō, rescrībere}
\textbf{rescrībō, rescrībere} — \textquotedblleft{}to write back, reply\textquotedblright{}.

\begin{quote}
Epistulam tuam accēpī et statim rescrībō.
\end{quote}

\begin{quote}
\emph{I have received your letter and I am writing back at once.}
\end{quote}

\emph{rescrībō, rescrībere, rescrīpsī, rescrīptum}. An emperor's written answer was a \emph{rescrīptum}.

\begin{cousinweb}
\textbf{rescrībō} is built on \emph{re-} \textquotedblleft{}back\textquotedblright{} and \emph{scrībō} \textquotedblleft{}to write\textquotedblright{} (Lewis \& Short). English \emph{rescript} comes from its participle.
\end{cousinweb}

\begin{grammarlens}[title={What you've built}]
One more word for telling, explaining and describing at B1.
\end{grammarlens}

\subsection*{Your turn}
\begin{itemize}
\raggedright
  \item \emph{Say it:} \emph{rescrībō, rescrībere}
  \item \emph{Say it:} \emph{rescrībō, rescrībere}, once more
  \item \emph{Say it:} a sentence of your own with \emph{rescrībō, rescrībere}
  \item \emph{Recall:} say the Latin for parchment, then the Latin for writing tablets, then say \emph{rescrībō, rescrībere} again
\end{itemize}

\begin{tcolorbox}[breakable,colback=teal!4,colframe=teal!35!black,title={Before you move on}]
What does rescrībō, rescrībere mean? (\textquotedblleft{}To write back, reply\textquotedblright{}.) Say it once more.
\end{tcolorbox}
